% !PRINTED 176
\weberpage{176}

If, with the help of the linear equation $b_1 = 0$, we
express one of the five quantities $\alpha$, say $\alpha_0$, in terms of the others, then we can regard the four remaining ones $\alpha_1, \alpha_2, \alpha_3, \alpha_4$ as homogeneous coordinates of a point in space. The equation
$b_2 = 0$ then represents a surface of the second order, $b_3 = 0$
a surface of the third order, $b_4 = 0$ a surface of the fourth order, $b_5 = 0$ a surface of the fifth order.

If one wished, then, to make $b_2$, $b_3$, $b_4$ vanish simultaneously and
thereby reduce the given equation to a pure one, one
would have to determine one of the 24 points of intersection of three surfaces
of the second, third, and fourth order, which depends on an equation of degree $24$.

Instead of proceeding in this way, one seeks to determine a straight line on the surface
$b_2 = 0$. Of such straight lines
there are on every surface of the second order one or two
families (real or imaginary), and one of these lines can be determined
\gesperrt{by square roots}, as will be carried out further below. The determination of a point of intersection of one
of these straight lines with the surface $b_3 = 0$ then leads to
an equation of the third degree, and if, by determining a common factor of the $\alpha$, one further makes the coefficient
$b_4 = 1$, the equation for $y$ takes the form
\weq{6}{y^5 + y + b = 0,}
whereby $y$ is determined as a function of \gesperrt{one} variable $b$.

This form of the equation is usually named after the English
mathematician \gesperrt{Jerrard}, who made it known in the year 1834.
But it had already been published much earlier (1786) by \gesperrt{E. J. Bring}
in Lund\footnote{Cf. F. Klein, Vorlesungen über das Ikosaëder, p. 143.}. We therefore call it the \gesperrt{Bring-Jerrard} form.

Just as well one could make $b_4 = 0$ by an equation of the fourth degree,
and would obtain a second normal form of the equation
of the fifth degree:
\weq{7}{y^5 + y^2 + b = 0,}
which has the same justification as the \gesperrt{Bring-Jerrard}
form.

For the solution of the equation of the fifth degree, to be sure, nothing
is gained by these considerations. Their usefulness consists
%CONT>>
